\documentclass[11pt,letterpaper]{article}

% ======= PACKAGES =======
\usepackage[utf8]{inputenc}
\usepackage{geometry}
\geometry{margin=1in}
\usepackage{setspace}
\setstretch{1.5}
\usepackage{fancyhdr}
\usepackage{xcolor}
\usepackage{titlesec}
\usepackage{lastpage}
\usepackage{hyperref}
\usepackage{mathpazo}
\usepackage{graphicx}
\usepackage{subcaption}
\linespread{1.5}

% ======= HEADER & FOOTER =======
\pagestyle{fancy}
\fancyhf{}

% Header (faded + smaller)
\fancyhead[L]{\footnotesize\textcolor{gray!90}{\textit{Ph.D. Abstract}}}
\fancyhead[R]{\footnotesize\textcolor{gray!90}{\textit{Mostafa Rezaali}}}

% Footer (faded)
\fancyfoot[L]{\footnotesize\textcolor{gray!90}{\textit{NSF AGS-PRF -- University of Florida}}}
\fancyfoot[R]{\footnotesize\textcolor{gray!90}{Page \thepage}}

\renewcommand{\headrulewidth}{0.4pt}
\renewcommand{\footrulewidth}{0.4pt}

% ======= PARAGRAPH STYLE =======
\setlength{\parindent}{15pt}
\setlength{\parskip}{6pt}
\titlespacing*{\section}{0pt}{0.45em}{0.15em}
\titlespacing*{\subsection}{0pt}{0.35em}{0.1em}

% ======= DOCUMENT BEGINS =======
\begin{document}

% ======= TITLE BLOCK =======
\begin{center}
    \vspace*{-1.5cm}
    {\Large\bfseries Postdoctoral Fellowship Application} \\[0.2cm]
    {\small\textit{NSF Atmospheric and Geospace Sciences Postdoctoral Research Fellowship}} \\[0.05cm]
    {\small\textit{Department of Geography, University of Florida}} \\[0.1cm]
    \rule{0.6\textwidth}{0.4pt}
\end{center}
\vspace{-0.5cm}

\noindent
Mostafa Rezaali \\
Ph.D. Candidate, Climate Science (Geography) \\
University of Florida \\
mostafarezaali@gmail.com \quad $\mid$ \quad (352) 709-5350 \\[0.3cm]
{\large\bfseries Ph.D. Abstract}\\[-0.15cm]
\rule{\textwidth}{0.4pt}

My doctoral research in Climate Science and Geography at the University of Florida focuses on artificial intelligence for extreme weather, especially heat waves and flash drought. The work develops and evaluates deep learning methods for spatiotemporal climate prediction, with attention to large-scale atmospheric drivers, local land-surface feedbacks, and forecast verification.

A major component is MeshFlowNet, a GraphCast-style graph neural network on an icosahedral mesh for direct 15-day prediction of CONUS daily maximum temperature anomalies. The model combines local PRISM temperature history, ERA5 atmospheric and land-surface variables, teleconnection indices, global circulation fields, topography, seasonal radiation, and land-mask information. My dissertation research uses cross-validated hindcasts, temporal anomaly correlation, uncertainty metrics, and baseline comparisons to determine where learned forecast models add value over climatology and persistence.

The broader goal is to improve physically interpretable and uncertainty-aware climate-risk prediction for heat, water, energy, agriculture, and public-health planning. This doctoral work provides the technical and scientific foundation for the proposed NSF AGS-PRF project on physics-informed conditional flow matching for probabilistic prediction of heat waves and flash-drought-relevant risk.
\end{document}
